\subsection{PROPERTIES PECULIAR TO CLOSED MAPPINGS}

PROPOSITION 9. Let X, X' be two topological spaces. A necessary and sufficient condition for a mapping \(f: \mathbf{X} \to \mathbf{X}'\) to be continuous and closed is that \(f(\overline{\mathbf{A}}) = \overline{f(\mathbf{A})}\) for every subset A of X.

The condition is sufficient, for it obviously implies that \(f\) is closed, and it also implies that \(f\) is continuous by reason of § 2, no. 1, Theorem 1. Conversely, if \(f\) is continuous and closed, we have \(f(\mathbf{A}) \subset f(\overline{\mathbf{A}}) \subset \overline{f(\mathbf{A})}\) by § 2, no. 1, Theorem 1; also \(f(\overline{\mathbf{A}})\) is closed in \(\mathbf{X}'\) by hypothesis; hence \(f(\overline{\mathbf{A}}) = \overline{f(\mathbf{A})}\) .

PROPOSITION 10. Let R be an equivalence relation on a topological space X. Then R is closed if and only if every equivalence class M mod R possesses a fundamental system of neighbourhoods which are saturated with respect to R.

Suppose R is closed, and let U be an arbitrary open neighbourhood of M; since \(F = \complement U\) is closed in X, the saturation S of F with respect to R is closed in X. Since M is saturated with respect to R, we have \(M \cap S = \varnothing\) , and thus \(V = \complement S\) is an open neighbourhood of M, saturated with respect to R and contained in U.

To prove the converse, let F be any closed subset of X. Let T be the saturation of F with respect to R, let x be a point of \(\complement T\), and let M be the equivalence class of x; then \(M \cap T = \varnothing\) and a fortiori \(M \cap F = \varnothing\) , so that U = \(\complement F\) is a neighbourhood of M. Hence there is a neighbourhood \(V \subset U\) of M such that V is saturated with respect to R; V does not meet F, hence does not meet T, so that \(\complement T\) is a neighbourhood of M and therefore of x. This shows that \(\complement T\) is open (§ 1, no. 2, Proposition 1), i.e.~T is closed.

Remark. Proposition 10 implies the following: if R is closed and if \(\varphi\) denotes the canonical mapping \(X \to X/R\) , then for each \(x \in X\) and each neighbourhood U of the equivalence class of \(x\) in X, \(\varphi(U)\) is a neighbourhood of \(\varphi(x)\) in X/R. It should be carefully noticed that this statement by no means implies that for every neighbourhood V of \(x\) , \(\varphi(V)\) is a neighbourhood of \(\varphi(x)\) ; in other words (no. 3, Proposition 5) a closed equivalence relation is not necessarily open (Exercise 2). Conversely, an open equivalence relation is not necessarily closed (no. 1, Example 4);

\pandocbounded{\includegraphics[keepaspectratio]{images/part001_bfecd507.jpg}}

for if U is a neighbourhood in X of an equivalence class M, then for each \(x \in M\) and each neighbourhood \(V \subset U\) of x, the saturation of V is certainly a neighbourhood of M in X, but this neighbourhood is not necessarily contained in U.

Finally, there are equivalence relations other than equality which are both open and closed (Exercise 3) and equivalence relations which are neither open nor closed ( \(§\ 8\) , Exercise 10).

